\documentclass[12pt,a4paper]{article}

% ---- Packages ----
\usepackage[margin=1in]{geometry}
\usepackage{graphicx}
\usepackage{listings}
\usepackage{xcolor}
\usepackage{hyperref}
\usepackage{amsmath}
\usepackage{booktabs}
\usepackage{fancyhdr}
\usepackage{titlesec}
\usepackage{float}
\usepackage{caption}

% ---- Colors ----
\definecolor{codegreen}{rgb}{0,0.6,0}
\definecolor{codepurple}{rgb}{0.58,0,0.82}
\definecolor{codegray}{rgb}{0.5,0.5,0.5}
\definecolor{backcolour}{rgb}{0.95,0.95,0.95}
\definecolor{darkblue}{rgb}{0.0,0.0,0.6}

% ---- Code Listings Style ----
\lstdefinestyle{verilogstyle}{
    backgroundcolor=\color{backcolour},
    commentstyle=\color{codegreen},
    keywordstyle=\color{darkblue}\bfseries,
    numberstyle=\tiny\color{codegray},
    stringstyle=\color{codepurple},
    basicstyle=\ttfamily\scriptsize,
    breaklines=true,
    numbers=left,
    numbersep=5pt,
    frame=single,
    tabsize=4,
    captionpos=b,
    morekeywords={module, endmodule, input, output, wire, reg, always, posedge, negedge, assign, begin, end, if, else, parameter, localparam, integer, for}
}
\lstset{style=verilogstyle}

% ---- Header / Footer ----
\pagestyle{fancy}
\fancyhf{}
\rhead{Assignment 10 -- Asynchronous FIFO}
\lhead{Digital IC Design}
\rfoot{Page \thepage}

% ---- Title Format ----
\titleformat{\section}{\Large\bfseries\color{darkblue}}{\thesection}{1em}{}
\titleformat{\subsection}{\large\bfseries\color{darkblue!80}}{\thesubsection}{1em}{}
\titleformat{\subsubsection}{\normalsize\bfseries\color{darkblue!60}}{\thesubsubsection}{1em}{}

% ---- Hyperref ----
\hypersetup{
    colorlinks=true,
    linkcolor=darkblue,
    urlcolor=blue,
    citecolor=darkblue
}

\begin{document}

% ===========================================================================
% TITLE PAGE
% ===========================================================================
\begin{titlepage}
\centering
\vspace*{2cm}
{\Huge\bfseries Asynchronous FIFO\\(ASYNC\_FIFO)\par}
\vspace{1cm}
{\Large Assignment 10 -- Digital IC Design Course\par}
\vspace{2cm}
{\large\textbf{Submitted by:} Ziad Mostafa Abdelaziz\par}
\vspace{0.5cm}
{\large\today\par}
\vfill
\end{titlepage}

\tableofcontents
\newpage

% ===========================================================================
% 1. INTRODUCTION & FIFO DEPTH CALCULATION
% ===========================================================================
\section{Introduction and Depth Calculation}

An \textbf{Asynchronous FIFO} is a dual-port memory buffer used to transfer data between two independent, asynchronous clock domains: a write domain (\texttt{W\_CLK}) and a read domain (\texttt{R\_CLK}). To prevent metastability and data incoherence across clock domain crossings (CDC), Gray-coded pointers are passed through double-flop synchronizers (\texttt{DF\_SYNC}).

\subsection{Mathematical FIFO Buffer Depth Calculation}

To prevent data overflow under burst conditions, the minimum required FIFO buffer depth is calculated based on the write and read frequency parameters given in the assignment specification:

\begin{itemize}
    \item \textbf{Write Clock Frequency ($f_w$):} $100\text{ MHz} \implies T_w = 10\text{ ns}$
    \item \textbf{Read Clock Frequency ($f_r$):} $40\text{ MHz} \implies T_r = 25\text{ ns}$
    \item \textbf{Burst Write Size ($B$):} $9\text{ Bytes}$ (written continuously at 100 MHz)
\end{itemize}

\subsubsection{Derivation Steps}

\begin{enumerate}
    \item \textbf{Burst Duration ($t_{\text{burst}}$):}
    \[
        t_{\text{burst}} = B \times T_w = 9 \times 10\text{ ns} = 90\text{ ns}
    \]

    \item \textbf{Items Read During Burst ($N_{\text{read}}$):}
    \[
        N_{\text{read}} = \left\lfloor \frac{t_{\text{burst}}}{T_r} \right\rfloor = \left\lfloor \frac{90\text{ ns}}{25\text{ ns}} \right\rfloor = 3\text{ items}
    \]

    \item \textbf{Minimum Buffer Capacity (Unread Items):}
    \[
        \text{Depth}_{\text{min}} = B - N_{\text{read}} = 9 - 3 = 6\text{ items}
    \]

    \item \textbf{Hardware Power-of-2 Depth Allocation:}
    In hardware, FIFO depth must be a power of 2 ($2^m$) to enable seamless binary-to-Gray pointer wrapping:
    \[
        \text{Hardware Depth} = 2^{\lceil \log_2(6) \rceil} = 8\text{ entries} \quad (\texttt{ADDR\_SIZE} = 3)
    \]
\end{enumerate}

\textbf{Conclusion:} A hardware FIFO depth of \textbf{8 entries} ($\texttt{ADDR\_SIZE} = 3$) guarantees zero overflow during the 9-byte burst write at 100 MHz with continuous reading at 40 MHz.

% ===========================================================================
% 2. BLOCK INTERFACE
% ===========================================================================
\section{Block Interface}

\subsection{Parameters}
\begin{table}[H]
\centering
\begin{tabular}{@{}llc@{}}
\toprule
\textbf{Parameter Name} & \textbf{Description} & \textbf{Default Value} \\
\midrule
\texttt{DATA\_WIDTH} & Width of data bus & 8 bits \\
\texttt{ADDR\_SIZE}  & Memory address size (Depth = $2^{\texttt{ADDR\_SIZE}}$) & 3 (Depth = 8) \\
\bottomrule
\end{tabular}
\caption{Design Parameters}
\end{table}

\subsection{Ports Description}
\begin{table}[H]
\centering
\begin{tabular}{@{}llc@{}}
\toprule
\textbf{Signal Name} & \textbf{Description} & \textbf{Width} \\
\midrule
\texttt{W\_CLK}   & Source domain write clock               & 1 \\
\texttt{W\_RST}   & Source domain active-low async reset    & 1 \\
\texttt{W\_INC}   & Write operation enable                  & 1 \\
\texttt{R\_CLK}   & Destination domain read clock          & 1 \\
\texttt{R\_RST}   & Destination domain active-low async reset & 1 \\
\texttt{R\_INC}   & Read operation enable                   & 1 \\
\texttt{WR\_DATA} & Write Data Bus                          & Parameterized (8 bits) \\
\texttt{RD\_DATA} & Read Data Bus                           & Parameterized (8 bits) \\
\texttt{FULL}     & FIFO Buffer full flag                   & 1 \\
\texttt{EMPTY}    & FIFO Buffer empty flag                  & 1 \\
\bottomrule
\end{tabular}
\caption{Port Descriptions}
\end{table}

% ===========================================================================
% 3. ARCHITECTURE AND SUB-MODULE DESIGN
% ===========================================================================
\section{Architecture and Design}

The top-level \texttt{ASYNC\_FIFO} integrates five modular Verilog blocks:

\begin{enumerate}
    \item \textbf{\texttt{FIFO\_MEM\_CNTRL}:} Dual-port RAM array ($8 \times 8$). Writes occur synchronously on \texttt{W\_CLK} when gated write enable $\text{wclken} = \text{W\_INC} \land \overline{\text{FULL}}$ is high. Reads occur asynchronously from location \texttt{raddr}.
    \item \textbf{\texttt{DF\_SYNC}:} 2-stage D flip-flop synchronizer chain used to safely transfer Gray pointers across asynchronous clock boundaries without metastability.
    \item \textbf{\texttt{FIFO\_WR}:} Generates binary write address \texttt{waddr}, Gray write pointer \texttt{wptr}, and evaluates the Clifford Cummings full condition:
    \begin{align*}
        \text{FULL} = & (\text{wgray\_next}[3] \neq \text{wq2\_rptr}[3]) \land (\text{wgray\_next}[2] \neq \text{wq2\_rptr}[2]) \\
                      & \land (\text{wgray\_next}[1:0] == \text{wq2\_rptr}[1:0])
    \end{align*}
    \item \textbf{\texttt{FIFO\_RD}:} Generates binary read address \texttt{raddr}, Gray read pointer \texttt{rptr}, and evaluates the empty condition:
    \[
        \text{EMPTY} = (\text{rgray\_next} == \text{rq2\_wptr})
    \]
    \item \textbf{\texttt{ASYNC\_FIFO}:} Structural top-level module connecting all sub-modules.
\end{enumerate}

% ===========================================================================
% 4. COMPARISON ANALYSIS (OUR VERSION VS. OLD VERSION)
% ===========================================================================
\section{Comparative Analysis: Our Version vs. Old Implementation}

A detailed comparative analysis was conducted between our optimized architecture and the previous implementation (\texttt{old version/ASYNC\_FIFO}):

\begin{table}[H]
\centering
\begin{tabular}{@{}lll@{}}
\toprule
\textbf{Design Feature} & \textbf{Old Implementation} & \textbf{Our Optimized Version} \\
\midrule
\textbf{Synchronizer} & 2D array + combinational \texttt{for} loop & Clean vector D-FF chain (\texttt{DF\_SYNC.v}) \\
\textbf{FULL Flag}    & Combinational \texttt{assign} (glitch-prone) & Registered on \texttt{W\_CLK} (glitch-free) \\
\textbf{EMPTY Flag}   & Un-registered combinational compare & Registered on \texttt{R\_CLK} (glitch-free) \\
\textbf{Spec Naming}  & \texttt{bit\_sync.v} (non-compliant) & \texttt{DF\_SYNC.v} (100\% spec compliant) \\
\textbf{Lint Status}  & 2D loop synthesis warnings & \textbf{0 Errors, 0 Warnings, 0 Latches} \\
\bottomrule
\end{tabular}
\caption{Architectural Comparison Summary}
\end{table}

\subsection{Key Design Advantages of Our Version}
\begin{enumerate}
    \item \textbf{Glitch-Free Flag Generation:} By registering \texttt{FULL} on \texttt{W\_CLK} and \texttt{EMPTY} on \texttt{R\_CLK}, transient combinational glitches across asynchronous domain bit changes are completely eliminated.
    \item \textbf{Synthesis Efficiency:} Using dedicated 1D vector registers in \texttt{DF\_SYNC.v} instead of 2D arrays avoids synthesis logic bloat and satisfies strict ASIC static timing constraints (STA).
    \item \textbf{Strict Spec Alignment:} Follows exact module naming (\texttt{DF\_SYNC.v}) and power-of-2 pointer parameterization.
\end{enumerate}

% ===========================================================================
% 5. RTL VERILOG CODE
% ===========================================================================
\section{RTL Verilog Modules}

\lstinputlisting[language=Verilog, caption={FIFO\_MEM\_CNTRL.v -- Dual-Port RAM Memory}]{../rtl/FIFO_MEM_CNTRL.v}
\lstinputlisting[language=Verilog, caption={DF\_SYNC.v -- Double Flop Synchronizer}]{../rtl/DF_SYNC.v}
\lstinputlisting[language=Verilog, caption={FIFO\_WR.v -- Write Pointer \& Full Flag Generator}]{../rtl/FIFO_WR.v}
\lstinputlisting[language=Verilog, caption={FIFO\_RD.v -- Read Pointer \& Empty Flag Generator}]{../rtl/FIFO_RD.v}
\lstinputlisting[language=Verilog, caption={ASYNC\_FIFO.v -- Structural Top Module}]{../rtl/ASYNC_FIFO.v}

% ===========================================================================
% 6. TESTBENCH AND SIMULATION RESULTS
% ===========================================================================
\section{Testbench and Simulation Results}

\subsection{Testbench Implementation}
The self-checking testbench (\texttt{ASYNC\_FIFO\_tb.v}) executes 4 comprehensive test cases simulating $100\text{ MHz}$ write clock and $40\text{ MHz}$ read clock.

\lstinputlisting[language=Verilog, caption={ASYNC\_FIFO\_tb.v -- Self-Checking Testbench}]{../tb/ASYNC_FIFO_tb.v}

\subsection{ModelSim Console Transcript}

All test cases passed with 0 errors:

\begin{figure}[H]
    \centering
    \includegraphics[width=\textwidth]{log output.png}
    \caption{ModelSim Transcript Log -- All 4 Test Cases PASSED (0 Errors)}
\end{figure}

\subsection{Output Waveform}

\begin{figure}[H]
    \centering
    \includegraphics[width=\textwidth]{wave form output.png}
    \caption{ModelSim Output Waveform -- ASYNC\_FIFO Simulation}
\end{figure}

\subsection{Detailed Waveform and Signal Transition Analysis}

The output waveform captures the complete asynchronous operation of all four test cases across independent clock domains:

\subsubsection{Write Domain Signals (100 MHz)}
\begin{itemize}
    \item \textbf{\texttt{W\_CLK} (Cyan):} The source domain clock runs at $f_w = 100$ MHz ($T_w = 10$ ns).
    \item \textbf{\texttt{W\_RST} (Orange):} Active-low asynchronous reset for the write domain. When \texttt{W\_RST} = 0, \texttt{wbin}, \texttt{wptr}, and \texttt{wfull} clear to 0 immediately.
    \item \textbf{\texttt{W\_INC} (Yellow):} Write enable signal. Pulsed high during burst writes of 9 bytes (\texttt{0x11} to \texttt{0x99}).
    \item \textbf{\texttt{WR\_DATA} (White, Hex):} Multi-bit write data bus. Holds test data bytes driven on \texttt{W\_CLK} edges.
    \item \textbf{\texttt{FULL} (Red):} FIFO full flag. Remains LOW (0) during TC2 concurrent read/write operation. In TC3 (full test), \texttt{FULL} asserts to HIGH (1) exactly when the write pointer reaches 8 entries without concurrent reads, inhibiting further memory writes.
\end{itemize}

\subsubsection{Read Domain Signals (40 MHz)}
\begin{itemize}
    \item \textbf{\texttt{R\_CLK} (Cyan):} The destination domain clock runs at $f_r = 40$ MHz ($T_r = 25$ ns).
    \item \textbf{\texttt{R\_RST} (Orange):} Active-low asynchronous reset for the read domain. Clears \texttt{rbin}, \texttt{rptr}, and sets \texttt{EMPTY} = 1.
    \item \textbf{\texttt{R\_INC} (Yellow):} Read enable signal. Enabled continuously during concurrent read operations.
    \item \textbf{\texttt{RD\_DATA} (Green, Hex):} Read data output bus from dual-port memory. Accurately outputs stored bytes in exact FIFO sequence (\texttt{0x11}, \texttt{0x22}, \dots, \texttt{0x99}).
    \item \textbf{\texttt{EMPTY} (Lime):} FIFO empty flag. Asserts to HIGH (1) post-reset and when all entries are drained (TC4), preventing read underflow.
\end{itemize}

\subsubsection{Internal Pointers and Synchronizer Signals}
\begin{itemize}
    \item \textbf{\texttt{waddr} \& \texttt{raddr} (Magenta, Unsigned):} 3-bit binary addresses ($0 \dots 7$) indexing memory locations.
    \item \textbf{\texttt{wptr} \& \texttt{rptr} (Gold, Binary Gray):} 4-bit Gray-code pointers. Only 1 bit changes per increment, eliminating multi-bit CDC metastability.
    \item \textbf{\texttt{wq2\_rptr} (Purple):} Read Gray pointer synchronized into the \texttt{W\_CLK} domain via \texttt{sync\_r2w} (2 D-FF stages).
    \item \textbf{\texttt{rq2\_wptr} (Purple):} Write Gray pointer synchronized into the \texttt{R\_CLK} domain via \texttt{sync\_w2r} (2 D-FF stages).
\end{itemize}

% ===========================================================================
% 7. SPYGLASS LINT RESULTS
% ===========================================================================
\section{SpyGlass Lint Results}

The design was analyzed using Synopsys SpyGlass (version vL-2016.06) under the \texttt{rtl\_handoff} methodology across all standard linting policies (lint, morelint, clock-reset, erc, latch, simulation, starc).

\subsection{Design Read Step Output}

SpyGlass successfully elaborated \texttt{ASYNC\_FIFO.v} and all sub-modules. The message tree shows \textbf{Total: 2 messages, Displayed: 2, Waived: 0} -- both are informational \texttt{Design Read} messages (no errors or warnings).

\begin{figure}[H]
    \centering
    \includegraphics[width=\textwidth]{spy glass Design read step output.png}
    \caption{SpyGlass Design Read Step Output -- 2 Info Messages, 0 Errors}
\end{figure}

\subsection{Lint RTL Goal Run Output}

The Lint RTL goal run completed with a \textbf{clean lint result} (0 Errors, 0 Warnings, 0 Latches).

\begin{figure}[H]
    \centering
    \includegraphics[width=\textwidth]{lint rtl goal output (lint clean).png}
    \caption{SpyGlass Lint RTL Goal Run Output -- Clean Lint (0 Errors, 0 Warnings)}
\end{figure}

\subsection{SpyGlass Lint Summary Report (moresimple.rpt)}

\begin{lstlisting}[language={}, caption={SpyGlass moresimple.rpt -- Lint Summary Report}]
Total Number of Generated Messages :          2
Number of Waived Messages          :          0
Number of Reported Messages        :          2
Number of Overlimit Messages       :          0

MORESIMPLE REPORT:

############### BuiltIn -> RuleGroup=Design Read ###############
[1] DetectTopDesignUnits  Info  ASYNC_FIFO.v  Line 8
    Module ASYNC_FIFO is a top level design unit
[0] ElabSummary           Info
    Please refer file elab_summary.rpt for elab summary report
\end{lstlisting}

% ===========================================================================
% 8. CONCLUSION
% ===========================================================================
\section{Conclusion}

The Asynchronous FIFO (\texttt{ASYNC\_FIFO}) was successfully designed, modeled, and verified.
\begin{itemize}
    \item Mathematical buffer depth calculation verified that an 8-entry FIFO ($\texttt{ADDR\_SIZE}=3$) fully accommodates a 9-byte burst write at 100 MHz with 40 MHz concurrent reading without overflow.
    \item Gray-code pointer encoding and double-flop synchronizers (\texttt{DF\_SYNC}) eliminated CDC metastability.
    \item Registered full/empty flags eliminate combinational glitches and provide superior timing margin over the old un-registered implementation.
    \item All 4 simulation test cases passed with 0 errors in ModelSim.
    \item Clean SpyGlass lint report: 0 errors, 0 warnings, 0 latches.
\end{itemize}

\end{document}
